\chapter*{Introducción}\label{chapter:introduction}
\addcontentsline{toc}{chapter}{Introducción}

El problema de partición de números (\emph{Number Partitioning Problem}, NPP) consiste en distribuir un conjunto de $N$ enteros positivos $a_1, a_2, \ldots, a_N$ en dos subconjuntos de modo que las sumas de sus elementos difieran lo menos posible. De manera equivalente, se busca una asignación de signos $\vec{s} \in \{-1,+1\}^N$ que minimice la discrepancia
\begin{equation}
    E(\vec{s}) = \nnorm{\sum_{i=1}^{N} s_i a_i}.
    \label{eq:discrepancy}
\end{equation}
Una partición se denomina perfecta cuando alcanza la menor discrepancia que admite la instancia, es decir, $E = 0$ si la suma total es par y $E = 1$ en caso contrario. La versión de decisión del problema, que pregunta si existe una partición perfecta, figura entre los problemas NP-completos clásicos \parencite{garey1979}. Sin embargo, su carácter NP-completo es solo débil: el problema puede resolverse mediante programación dinámica en un tiempo polinomial no únicamente en $N$, sino también en la magnitud de los números de entrada \parencite{garey1979,mertens2001}. La dificultad de una instancia depende, por tanto, no solo de su cardinalidad sino también de la cantidad de bits $b$ necesaria para representar sus elementos.

Más allá de su interés teórico, el NPP posee aplicaciones prácticas en la planificación de tareas sobre multiprocesadores y en la minimización del tamaño y el retardo de circuitos VLSI \parencite{gent1996,mertens2001}, así como en la construcción de criptosistemas de clave pública basados en el problema de la mochila \parencite{merkle1978,mertens2006}, emparentado con el NPP. Su relevancia teórica se debe, además, a que es el único de los problemas básicos de la teoría de la NP-completitud formulados por \textcite{garey1979} que trata directamente con números, lo que lo convierte en elección habitual para demostrar la NP-completitud de otros problemas numéricos \parencite{gent1996,mertens2006}.

La dependencia respecto de $b$ se formaliza al considerar instancias aleatorias en las que cada $a_i$ se extrae de forma independiente y uniforme del conjunto $\{1, \ldots, 2^{b}\}$. En este ensemble el comportamiento del problema queda gobernado por el parámetro $\kappa = b/N$, que expresa la cantidad de bits por número en relación con el tamaño de la instancia. En el límite $N \to \infty$ existe un valor crítico $\kappa_c = 1$ que separa dos regímenes cualitativamente distintos \parencite{mertens1998,mertens2001,borgs2001}. Para $\kappa < \kappa_c$ casi todas las instancias poseen un número exponencial de particiones perfectas, de modo que un algoritmo de búsqueda suele localizar una de ellas con rapidez \parencite{mertens1998,mertens2001}. Para $\kappa > \kappa_c$, en cambio, casi ninguna instancia admite una partición perfecta y la solución óptima es única, salvo por la simetría global de signos \parencite{borgs2001}, y difícil de hallar \parencite{mertens1998,mertens2001}. Para tamaños finitos el valor crítico presenta correcciones, que \textcite{mertens1998} cuantifica mediante $\kappa_c(N) = 1 - \log_2(\pi N/6)/(2N)$ y a las que atribuye la estimación experimental $\kappa_c \approx 0{,}96$ obtenida por \textcite{gent1996} con instancias de hasta $N = 30$. Esta transición entre una fase fácil y una fase difícil, que refleja la pseudopolinomialidad del problema \parencite{mertens1998}, constituye el objeto de estudio de esta tesis. El caso de números reales distribuidos uniformemente en $[0,1)$, habitual en los análisis teóricos, corresponde al límite de precisión infinita en el que $\kappa \to \infty$ y el problema se sitúa por completo en la fase difícil \parencite{mertens2001}; su estudio está mas enfocado en el comportamiento de la energía y se dejan de perseguir particiones perfectas.

La idea de que la dificultad computacional típica de un problema NP-completo se concentra en torno a una transición de fase del ensemble aleatorio se ha documentado en numerosos problemas, entre ellos la satisfacibilidad booleana y el circuito hamiltoniano \parencite{gent1998,mertens2006}. Para el NPP se sostuvo inicialmente, a partir de un análisis de física estadística debido a Fu, que no existía transición alguna. \textcite{gent1996} refutaron esta conjetura: mediante una teoría \emph{annealed} definieron un parámetro de restricción $\kappa$ y localizaron experimentalmente la transición, trabajo que ampliaron en \textcite{gent1998}. Poco después, \textcite{mertens1998} estableció mediante un análisis de física estadística la existencia de la transición, con valor crítico $\kappa_c = 1$, y la vinculó con la pseudopolinomialidad del problema; este análisis se expone de forma más amplia en \textcite{mertens2001}, y \textcite{borgs2001} la establecieron con rigor matemático, demostrando que es de primer orden y determinando una ventana de escalado de ancho $\Theta(1/N)$ centrada en $\kappa_N = 1 - \log_2 N/(2N)$, en lo que sus autores describen como el primer análisis riguroso de escalado de tamaño finito en un problema NP-completo. \textcite{ferreira1998} abordaron el problema aprovechando su equivalencia con la búsqueda del estado base de un modelo de Ising de alcance infinito con interacciones aleatorias, y obtuvieron cotas para la discrepancia óptima media mediante la aproximación \emph{annealed}.

Más allá de las propiedades del ensemble, la literatura aporta evidencia sobre el comportamiento de los algoritmos en torno a la transición. \textcite{gent1998} mostraron que el esfuerzo de búsqueda de Complete Karmarkar--Karp (CKK) alcanza un máximo en torno a la transición y se mantiene elevado en la fase insoluble, y que el reescalado de tamaño finito de la calidad de las soluciones heurísticas requiere puntos críticos distintos para cada heurística, aproximadamente $0{,}40$ para la de Karmarkar--Karp y $0{,}15$ para la voraz --que asigna el mayor número libre a la partición de menor suma--, frente al valor próximo a $0{,}96$ de la existencia de particiones perfectas. \textcite{korf1998} observó, por su parte, que las instancias más difíciles para CKK aparecen allí donde la probabilidad de una partición perfecta es próxima a $0.5$. \textcite{mertens2006} indica que la propia heurística de Karmarkar--Karp solo genera particiones perfectas por debajo de un umbral que depende de $N$ y que parece tender a cero cuando $N \to \infty$, y que en la fase difícil los algoritmos completos basados en las heurísticas voraz y de diferenciación resultan equivalentes a la búsqueda aleatoria. \textcite{mertens2001} resume además que el tiempo de solución exacta crece exponencialmente con $N$ para $N \lesssim b$ y polinomialmente para $N \gg b$, con un máximo cuando $N \sim b$. La síntesis de estos resultados teóricos se encuentra en \textcite{mertens2001}, trabajo que constituye la base conceptual de esta tesis.

En paralelo se desarrollaron los algoritmos cuyo comportamiento se examina aquí. El método de diferenciación de \textcite{karmarkar1982} sustituye iterativamente los dos mayores números del conjunto por su diferencia, lo que da lugar a una heurística de coste $O(N \log N)$ \parencite{mertens2001,korf1998} que resulta notablemente más efectiva que la heurística voraz \parencite{korf1998,gent1998}. \textcite{korf1998} extendió esta idea hasta el algoritmo completo CKK, que explora sistemáticamente las dos formas de combinar cada par de números, mediante su diferencia o su suma, y que produce soluciones de calidad creciente a medida que se prolonga la búsqueda; según su autor, resuelve en la práctica instancias de tamaño arbitrario con hasta doce dígitos decimales de precisión. Por su parte, el recocido simulado, propuesto por \textcite{kirkpatrick1983} a partir de la analogía entre la optimización combinatoria y la mecánica estadística, y apoyado en el algoritmo de Metropolis y en los métodos de Monte Carlo \parencite{newman1999}, fue sometido a una evaluación experimental sistemática organizada en una serie de artículos \parencite{johnson1989,johnson1991}. En el segundo de ellos, que incluye el NPP, se concluyó que el recocido simulado no resulta competitivo frente al algoritmo de diferenciación salvo en instancias pequeñas y que, a igualdad de tiempo de cómputo, ni siquiera supera a múltiples ejecuciones de optimización local desde inicios aleatorios \parencite{johnson1991}. \textcite{ferreira1998} ofrecen una explicación de este fracaso: los movimientos de un solo espín alteran la energía en un orden $1/N$, mientras que la discrepancia óptima es del orden de $\sqrt{N}\,2^{-N}$. Un resultado reciente, de especial relevancia para este trabajo, es el de \textcite{angelini2026}, quienes muestran para el recocido simulado aplicado al problema aleatorio de $K$-satisfacibilidad que los umbrales algorítmicos dependen de la forma en que el tiempo de ejecución escala con el tamaño del sistema. De ello se desprende que, para un par problema--algoritmo dado, la frontera entre la región accesible y la inaccesible no es única, sino que depende del presupuesto de cómputo concedido al método.

La transición del NPP está bien establecida como propiedad del ensemble, pero su manifestación en el desempeño de algoritmos concretos depende de la naturaleza de cada uno. Cabe esperar que un algoritmo completo la refleje en el esfuerzo necesario para certificar la optimalidad, que una heurística constructiva de tiempo fijo lo haga únicamente en la calidad de las soluciones que produce y que un método estocástico lo haga, además, en su dependencia del tiempo de búsqueda. Resulta pertinente, por tanto, un análisis conjunto, bajo un mismo diseño experimental, que confronte métodos de distinta naturaleza y permita discernir qué rasgos de la transición son propios del problema y cuáles dependen del algoritmo. Un estudio de esta índole aporta evidencia empírica reproducible sobre una cuestión que, en términos generales, permanece abierta: la conexión entre las transiciones de fase del espacio de soluciones y la dificultad práctica que enfrentan los algoritmos no está aún claramente establecida \parencite{angelini2026}. Ofrece, además, un criterio para elegir el método más adecuado según la región del espacio de parámetros en que se halle la instancia.

La problemática central de la tesis puede formularse, por tanto, como la interrogante de si la transición de fase del NPP se refleja de manera observable en el desempeño de la heurística de Karmarkar--Karp, del algoritmo completo de Karmarkar--Karp basado en ramificación y poda y del recocido simulado y, de ser así, si ocurre en la misma región del parámetro $\kappa$ para todos ellos. De esta pregunta se derivan otras más específicas. Se requiere determinar a través de qué magnitudes se manifiesta el cambio de fase en cada algoritmo, ya sea el esfuerzo computacional, la calidad de la solución o la probabilidad de hallar una partición perfecta, y con qué métricas puede cuantificarse de forma comparable. Interesa asimismo establecer cómo varían la posición y la nitidez de la transición observada al aumentar $N$, lo que remite a su escalado de tamaño finito \parencite{borgs2001,gent1998}. En el recocido simulado, cabe indagar hasta qué punto el cambio de fase depende del número de iteraciones y del esquema de enfriamiento adoptados, a la luz de los resultados sobre la dependencia de los umbrales respecto del escalado temporal \parencite{angelini2026}.

El objetivo general de la tesis es analizar de forma empírica el comportamiento de estos tres algoritmos en las proximidades de la transición de fase del problema de partición de números, con el propósito de comprobar su existencia en cada uno de ellos y de valorar si se mantiene de manera consistente entre métodos. Para alcanzarlo se caracteriza la distribución aleatoria y su parámetro de control a partir de la literatura de referencia, se implementan los algoritmos y se definen las magnitudes que permiten observar la transición en cada caso. A continuación, se diseña y ejecuta una campaña experimental que recorre $\kappa$ a ambos lados de su valor crítico para varios tamaños de instancia con números enteros de $b$ bits, complementada con instancias de números reales uniformes en $[0,1)$, que corresponden al límite de precisión infinita. A partir de los resultados se estiman la ubicación de la región de transición correspondiente a cada método, se contrastan estas estimaciones entre sí y con el valor teórico $\kappa_c = 1$, y se interpretan las coincidencias y discrepancias observadas a la luz de los resultados teóricos disponibles, con vistas a valorar el alcance de la transición de fase como descriptor de la dificultad algorítmica del problema.

El documento se organiza en capítulos que acompañan el desarrollo lógico del trabajo. El Capítulo~\ref{chapter:background} presenta el marco teórico, que comprende la formulación del problema, el ensemble aleatorio, la teoría de la transición de fase y los fundamentos de los algoritmos estudiados. El Capítulo~\ref{chapter:design} describe el diseño experimental, esto es, la generación de instancias, la implementación de los métodos, las métricas y el protocolo de ejecución. El Capítulo~\ref{chapter:results} expone y discute los resultados obtenidos. Cierran el documento las Conclusiones y las Recomendaciones, estas últimas acompañadas de las líneas de trabajo futuro que se derivan del estudio.
